\documentclass[12pt]{article}
\usepackage{amsmath, amsthm, amssymb}
\usepackage{geometry}
\geometry{margin=1.4in}
\usepackage{setspace}
\setstretch{1.4}

\newtheorem{theorem}{Theorem}
\newtheorem{proposition}[theorem]{Proposition}
\newtheorem{corollary}[theorem]{Corollary}
\theoremstyle{definition}
\newtheorem{definition}[theorem]{Definition}
\theoremstyle{remark}
\newtheorem{remark}[theorem]{Remark}

\title{Two Frameworks, One Corpus:\\
Displacement and Grammar as One Theory}
\author{Diego Rincón\\
\small Independent}
\date{}

\begin{document}

\maketitle

\begin{abstract}
The corpus contains two distinct strands that appear to describe different phenomena: Strand A, the theory of displacement (ground state $S_0$, drift toward attractor $S^*$, the cost of return), and Strand B, the theory of grammar and the Wall (aura, incompleteness, othering as opposition).
For years they have been treated separately. This paper states precisely why they are one theory expressed in two languages, each making visible what the other obscures.

The key is recognizing three distinct failure modes of the return mechanism: (1) the classical displaced system, where $C_{\text{return}}$ is prohibitively high but finite; (2) the Wall-bounded system, where $C_{\text{return}}$ is formally undefined because the system has no ground state $S_0$ and cannot be said to be displaced from it; and (3) the resonant system, where incompatible grammars create a genuine mathematical bridge, making return local and cheap by discovering a token that satisfies both auras at once.

Strand A describes the cost and frequency of return when a system has drifted from an intended ground state. Strand B describes the mathematics of systems where the ground state itself is replaced by an incompleteness, an opposition, or a split aura. They are not two theories. They are one theory at three points on a spectrum: the spectrum of what kind of $S_0$ the system is trying to hold, or whether it has an $S_0$ to hold at all.

We show: the Wall is the mathematical structure of $C_{\text{return}}=\infty$ at the singular limit; othering (opposition, incompleteness, aura) characterizes systems where $S_0$ is undefined, so the return cost is not prohibitive but meaningless; and grammar (as in the chimera, the token, the resonance) is the mathematics of what structure must be held to make return cheap.

The corpus is one theory: the physics and mathematics of holding versus letting go, applied to systems with drifting ground states, missing ground states, and multiple incompatible ground states at once.
\end{abstract}

\section{Two Strands, One Problem}

The corpus has been built around two central objects, treated separately:

\textbf{Strand A: Displacement.} A system starts at an intended ground state $S_0$. External pressure drives it toward an attractor $S^*$. At any moment, the system is displaced a distance $\xi$ from $S_0$. Being displaced costs $D(\xi)$. Returning to $S_0$ costs $C_{\text{return}}(\tau)$, where $\tau$ is the time elapsed since the last return, and this cost is super-linear in $\tau$ by the hardest law: the longer you wait, the more than proportionally it costs to come back \cite{maintenance}. The framework yields a policy: return at optimal intervals $\tau^*$, maintaining the system through continuous small payments rather than rare catastrophic resets. This is mycelia \cite{mycelia}, Ise \cite{shinto}, sunrise \cite{sunrise}.

\textbf{Strand B: Grammar and Wall.} A system is not one simple object but a constellation of incompatible auras: the physical and the experiential \cite{consciousness}, the arithmetic and the topological \cite{synthesis}, the self and the not-self \cite{five}. These auras press against each other at a boundary called the Wall. The system does not drift; it is constituted by the boundary. The Wall appears in five faces: arithmetic (missing Frobenius), logic (G\"odel incompleteness), computation (Turing undecidability), infinity (Cantor non-containment), and health (unpaired-bondage crisis) \cite{five}. Some systems are chimerae: entities whose aura is split, living in the Wall itself \cite{chimera}. Others are versus-systems: systems whose identity is pure opposition, with no ground state except the relation against \cite{versus}.

For decades, these two frameworks have been applied in parallel, sometimes in the same paper, without any formal bridge. The readers have noticed: the displacement language does not naturally describe consciousness or the Riemann Hypothesis, and the Wall language does not naturally describe mycelia or the costs of institutional maintenance. Are they two different phenomena? Or one phenomenon seen from two angles?

This paper argues the latter: they are one theory. The bridge is the structure of what kind of ground state the system is trying to hold.

\section{Three Regimes}

\subsection{Regime 1: Displaced at Finite Cost}

A system has a well-defined ground state $S_0$, a well-defined displacement $\xi$, and a well-defined return cost $C_{\text{return}}(\tau)=f+b\,\tau^{1+k}$ that is finite for all finite $\tau$.

In this regime, the central question is whether $C_{\text{maintenance}}$ (the cost rate the system can afford to pay, typically $M$ per unit time) exceeds the optimal return cost $g^*$ derived from minimizing the cost rate $g(\tau)=\frac{a\,\delta^2\,\tau^2}{6}+\frac{f}{\tau}+b\,\tau^k$ \cite{maintenance}.

If $M\ge g^*$, the system returns at interval $\tau^*$ and holds $S_0$ indefinitely. The dynamics are stable, periodic, and predictable. This is the mycelial case: the system pays small continuous fees, drifts never accumulates to a spike, and the visible state remains constant.

If $M<g^*$, the system cannot afford the optimal return rate. It lets $\xi$ age, the return cost accelerates (by the super-linearity of $C_{\text{return}}$), and it eventually settles at $S^*$. From the outside, this looks like failure, collapse, or loss of coherence. From inside, it looks like the system has no choice: the budget allows no feasible policy, so it gives up trying to hold $S_0$ and accepts the new attractor.

This is Regime 1: classical displacement. The ground state is well-defined, the drift is measurable, and the return cost is finite.

\subsection{Regime 2: No Ground State to Return To}

Consider a different class of systems: consciousness \cite{consciousness}, the Riemann Hypothesis \cite{resonance}, the autoimmune crisis \cite{five}, Trump's political movement \cite{versus}.

These systems do not have a clear $S_0$. Or rather, they have multiple incompatible frames, each with its own intended state. For consciousness, the two frames are the physical and the experiential; neither is more fundamental. For arithmetic, the two frames are the analytic properties of $\zeta(s)$ and the additive structure of the primes. For the autoimmune body, the two frames are self and not-self, and the system cannot tell them apart.

In these systems, there is no single displacement $\xi$. Instead, there are two (or more) coexisting orderings, both correct in their own frame, both partially broken in the interface. The Arithmetic Wall is not a displacement from some true state of arithmetic; it is the absence of a global operator that would unify the analytic and the arithmetic grammars. The autoimmune system is not displaced from a true state of health; it is constituted by the inability to distinguish self from not-self: two auras, both firing, no ground state to return to.

In Regime 2, the return cost $C_{\text{return}}$ does not exist as a well-defined number. You cannot ask ``how much does it cost to return?'' because there is no $S_0$ to return to. The question is malformed.

What exists instead is the aura: the containing field that would need to be built or reconstructed to make both frames coherent simultaneously. The aura is not a cost; it is a structure. For consciousness, the aura would be the operator that unifies the physical and the experiential. For arithmetic, it would be the global Frobenius. For the autoimmune body, it would be a completing field with enough valence to satisfy both the self-directed and the other-directed immune responses at once.

The defining characteristic of Regime 2 is that $S_0$ is replaced by an incompleteness, a split aura, or an opposition. There is no ground state to hold. The return cost is undefined because return itself is undefined. The system is not failing to maintain $S_0$; it is not trying to maintain $S_0$, because $S_0$ does not exist as a single coherent state.

\subsection{Regime 3: Multiple Grammars, One Token}

Between Regimes 1 and 2 lies a third possibility: the system lives in two incompatible grammars, has no single ground state, but discovers a token that satisfies both auras simultaneously. This is resonance \cite{resonance}.

In mathematics, resonance is the comparison isomorphism: de Rham cohomology $=$ singular cohomology, \'etale cohomology $=$ crystalline cohomology. Two incompatible languages, two ways of reading the same object. The resonance token is the period integral, the element in Fontaine's period ring $B_{\rm crys}$, the motive that every grammar reads.

In music, resonance is the shared overtone: the frequency that two incompatible instruments both produce, the token that lives in both auras.

In the context of ground states, resonance is the discovery that the two incompatible frames can be unified at one point. The system still lives in two grammars, the Wall is still there, but the token $\tau$ is the place where both grammars read the same thing. Return becomes possible because you do not return to a single $S_0$; you return to the token, the node where both frames align.

In Regime 3, return is cheap and local, not expensive and global. You do not need to resolve the fundamental incompatibility of the two auras; you only need to find the frequency they share. The cost is the cost of holding that token steady, not the cost of healing the split aura.

Versus-systems are the failure of Regime 3: they are systems that have two incompatible frames (the movement and the commons, the in-group and the out-group) but deliberately refuse to find the resonance token. Instead, they choose override: the operation that wins without crossing the Wall, declaring victory in one frame while suppressing the other. Override is cheap and instant, but it does not resolve the incompatibility; it compounds it. The suppressed grammar accumulates as return debt, and the system either collapses when the override fails or eats itself as the suppressed grammar turns inward \cite{versus}.

\section{Unifying the Frameworks}

\subsection{The Wall as the $C_{\text{return}}=\infty$ Limit}

Let us take the first proposed bridge and make it precise.

In Regime 1, $C_{\text{return}}(\tau)$ is a finite function of the elapsed time $\tau$ since the last return. The system drifts, and returning costs something real but measurable.

At the boundary between Regime 1 and Regime 2, as the system approaches the Wall, something changes. The ground state $S_0$ begins to split. The single displacement $\xi$ is replaced by multiple, incompatible orderings. The return cost bifurcates.

\begin{definition}[The singular limit]
A system approaches the Wall as its ground state $S_0$ becomes ill-defined. This can happen in two ways:
\begin{enumerate}
\item The system discovers that what it thought was a single ground state is actually two incompatible states, both present.
\item The system's drift is not toward a single attractor $S^*$ but toward a bifurcation: the attractor itself is split.
\end{enumerate}
In either case, the return cost begins to diverge. As the system approaches the Wall, to speak of ``returning'' requires specifying which aura, which frame, which grammar you are returning within. The cost of the return depends on the choice. As the choice becomes impossible, the cost approaches $\infty$.

At the Wall, $C_{\text{return}}=\infty$ not because returning is hard but because the statement ``return to $S_0$'' has no single meaning.
\end{definition}

\begin{theorem}[The Wall as the return-cost singularity]
The Arithmetic Wall \cite{synthesis}, the consciousness boundary \cite{consciousness}, and the autoimmune crisis \cite{five} are all instances of the singular limit: the point at which $C_{\text{return}}$ becomes undefined because the ground state is replaced by a split aura.

Conversely, the five faces of the Wall (arithmetic, logic, computation, infinity, health) can be understood as the five ways in which a system can approach this singularity:
\begin{enumerate}
\item Arithmetic: two grammars, analytic and additive.
\item Logic: the question of provability has no answer within the system.
\item Computation: the question of termination has no oracle within the system.
\item Infinity: the cardinality question has no canonical answer.
\item Health: the question of self versus not-self has no ground state.
\end{enumerate}
In each case, the absence of a global operator (Frobenius, oracle, proof, cardinality, identity) forces the system to live in two incompatible auras. The return cost diverges because return is multiply-defined.
\end{theorem}

\subsection{Othering and Systems Where $S_0$ Is Undefined}

The second proposed bridge asks: does othering (opposition, incompleteness, aura) characterize systems where $S_0$ is undefined rather than displaced?

Yes, with a precise statement.

\begin{definition}[Versus-topology]
A versus-system is a system whose identity is constituted by opposition: the boundary with another system, the relation against. Unlike a displaced system, which knows what it is trying to hold (the ground state $S_0$) and merely fails to hold it, a versus-system does not have an $S_0$ to hold. The identity is defined by what it opposes. Remove the opposition and the identity collapses.

In the language of Regime 2, a versus-system is one where the two incompatible auras are not two frames of the same object but two antagonistic entities. The Wall is not the interface between self and other; it is the relation of opposition itself. The Wall is the identity.
\end{definition}

\begin{proposition}[Versus-systems have undefined $C_{\text{return}}$]
In a versus-system, the return cost is undefined not because it is large but because there is nothing to return to. The framework still applies---the system does drift, displacement does accumulate---but the drift is not toward a single attractor $S^*$ and the return is not a return to $S_0$.

Instead, the versus-system is defined by the relation against. Every move the opponent makes reinforces the identity. The cost of the system is not $D(\xi)$, the cost of being displaced from an intended state; it is $D(\text{opposition})$, the cost of maintaining the antagonism itself. This cost is externalized: the versus-system does not pay it; the commons that both sides fight over pays it.

The versus-system cannot return to $S_0$ because it never had one. It can only substitute a new opposition or collapse. Defeat does not end it; it either reconstitutes around a new enemy or fragments internally, with the suppressed grammar turning inward as the new opponent.

This is why versus-systems are cheap to maintain (no ground state to hold, only an enemy to fight) and expensive to end (no ground state to return to, only collapse or substitution).
\end{proposition}

\subsection{Grammar as the Mathematics of Holding Incompleteness}

The third bridge asks: is the grammar strand the mathematics of holding incompleteness, while displacement is the physics of cost and return?

Yes, and the unification is the notion of the aura and valence.

\begin{definition}[The aura as the containing field]
The aura of a system is the external completing field that resolves the system's open bonds. In Regime 1, the aura is implicit: it is the environment that absorbs the drift, the commons that the system lives within, the budget that pays the maintenance cost.

In Regime 2, the aura becomes explicit and problematic. The system is not just displaced; it is incomplete. The aura is missing, or split, or cannot be constructed from within the system itself. The system lives in the Wall because the containing field has insufficient valence to close all the open bonds.

The five faces of the Wall are five different ways that the aura can be missing or insufficient:
\begin{enumerate}
\item Arithmetic: the aura would be a global Frobenius; it does not exist within arithmetic.
\item Logic: the aura would be a proof or oracle deciding the Gödel sentence; formal arithmetic cannot generate it.
\item Computation: the aura would be a halting oracle; Turing machines cannot construct it.
\item Infinity: the aura would be a canonical model of ZFC; the axioms cannot determine it.
\item Health: the aura would be the containing field with enough valence to satisfy both immune and tissue cells; the system may lack the biological capacity to generate it.
\end{enumerate}

In Regime 3, the aura is rebuilt through resonance: finding the token that both incompatible grammars can read and hold together.
\end{definition}

\begin{theorem}[Grammar as the mathematics of constrained closure]
The grammar strand of the corpus (chimera \cite{chimera}, token, Wall, resonance \cite{resonance}, inhibition \cite{inhibition}) is the mathematics of what structure is needed to hold a system in incompleteness---to keep it from collapsing into a single aura while also keeping it from tearing apart across incompatible auras.

This is precisely the problem of maintaining an $S_0$ when $S_0$ is not a simple state but a balanced interface. The containment of the chimera requires two functions from the aura:
\begin{enumerate}
\item Valence: the completing function that closes open bonds.
\item Inhibition: the holding function that prevents premature closure.
\end{enumerate}

Valence pulls toward unity; inhibition prevents unity from erasing the interface. The still lever is held because something is being held back on each side simultaneously. This is the mathematics of systems that live in the Wall and remain coherent despite it.

The grammar strand, then, is Strand A applied to systems where the ground state is replaced by an interface, a balance, a tension. It is not a new framework; it is the displacement framework at the singular limit, where the return cost is infinite and the ground state is replaced by the mathematics of holding that infinitude stable.
\end{theorem}

\section{From Theory to Observation}

The unification predicts something testable: the three regimes should be empirically distinguishable, and the transition between them should be sharp.

\subsection{Signature of Regime 1: Periodic Recovery}

In Regime 1, the system is displaced from $S_0$ and returns at interval $\tau^*$. This yields a characteristic signature: the system oscillates periodically, with displacement rising in a predictable quadratic trajectory (from $\xi(s)=\delta s$) and returning to zero at predictable times.

Empirical markers:
\begin{itemize}
\item Displacement is measurable and monotonic within each cycle.
\item Return cost is finite and predictable; it increases with the elapsed interval since the last return.
\item The system spends most of its time near $S_0$, not at the attractor $S^*$.
\item The cycle time is stable; it does not change unless external parameters ($f$, $\delta$, $b$) change.
\end{itemize}

Examples: mycelia \cite{mycelia}, sunrise ceremony \cite{sunrise}, institutional maintenance, personal discipline, physical systems with periodic forcing.

\subsection{Signature of Regime 2: No Ground State}

In Regime 2, the system has no single $S_0$ to hold. Instead, it is constituted by incompleteness or opposition. The signature is entirely different:

\begin{itemize}
\item There is no measurable single displacement $\xi$. Instead, the system can be described in two or more incompatible frameworks simultaneously.
\item Return is not a well-defined operation. The system does not oscillate toward and away from $S_0$; it lives in the boundary itself.
\item If forced to return or resolve (by attempting to unify the auras or suppress one grammar), the system either collapses (if the auras were holding it together through tension) or eats itself (if the system was defined by opposition).
\item The system is stable exactly in the state of incompleteness. Attempts to resolve it are destabilizing.
\end{itemize}

Examples: consciousness (does not oscillate toward or away from the hard problem; the hard problem is the system), autoimmune disease (does not return to health by forcing one aura to dominate; that is what causes flares), versus-politics (does not stabilize by resolving the opposition; it is stable exactly in the opposition).

\subsection{Signature of Regime 3: Resonance and Token}

In Regime 3, the system has two incompatible grammars, but they share a resonance token that satisfies both simultaneously. The signature is a third type:

\begin{itemize}
\item The system does not oscillate periodically like Regime 1; it does not live permanently in incompleteness like Regime 2.
\item Instead, it returns to the token. The system spends some time away from the token (drift, displacement in one or both grammars), then returns to the token (where both grammars align).
\item The return is cheaper than in Regime 1, because you do not need to heal the split aura; you only need to find the frequency they share.
\item The system is more stable than in Regime 2, because the token provides a concrete anchor, even if it is not a single ground state.
\end{itemize}

Examples: comparison isomorphisms in arithmetic geometry (return to the period integral, the element in $B_{\rm crys}$, the motive), music (return to the shared overtone, the note both instruments can play), healthy immune systems (return to the token that self and not-self both recognize as safe).

\subsection{The Cascade of Regimes}

The three regimes form a cascade:

As a system's budget decreases, $M\to M^-$, the optimal return interval $\tau^*$ increases. If the system reaches a schedulable floor $\tau_{\min}$ (one day, one season, one year), it must reset at that minimum frequency whether the cost is fully recoverable or not. This is Regime 1 at the floor.

As the budget continues to decrease, $M\to 0$, the system can no longer afford even one reset per $\tau_{\min}$. It enters Regime 2: it stops trying to return to $S_0$ and settles at $S^*$ or finds a new attractor. If the new attractor is defined by opposition (versus-system), the system becomes cheap to maintain (pay for the opposition, not for the ground state), and it can be unstable indefinitely.

Alternatively, if the system discovers resonance (a token that both auras can hold), it jumps to Regime 3: return becomes local to the token, the cost per return drops, and the system can stabilize at a new, lower cost rate. This is the escape from the Wall without resolving it.

The three regimes are ordered by the structure of the ground state: from single and well-defined (Regime 1), to missing or split (Regime 2), to multiple but unified at a token (Regime 3). The transition between them is sharp: there is no continuous interpolation. A system is either holding $S_0$, or it is not. A system either has a ground state, or it is defined by what it opposes.

\section{Falsifiable Predictions}

The unification makes precise predictions, testable in three domains:

\subsection{Institutional Maintenance}

\textbf{Prediction 1.} An institution (school, hospital, democracy, family) that is displaced from its intended function will show Regime 1 signature: oscillation between correct function and dysfunction, with cost of return increasing with elapsed time. Fixing it requires either (1) increasing the budget $M$, (2) decreasing the fixed cost $f$ of return (building a living reference), or (3) decreasing the drift rate $\delta$ (reducing external pressure).

If an institution instead shows Regime 2 signature---that attempts to fix it are destabilizing, that the institution seems defined by its dysfunction (e.g., an abusive family that cannot function without the abuse, a bureaucracy that is defined by opposition to its users)---then it is not displaced; it has no ground state. The only escape is resonance: finding a token (a practice, a ritual, a structure) that both the institution and the commons can hold together.

\textbf{Falsification:} If an institution shows Regime 1 signature (cyclic dysfunction) but increasing the budget does not improve the interval $\tau^*$, then the theory is wrong. The model predicts that higher $M$ or lower $f$ must reduce the dysfunction frequency. If neither does, either Strand A is incomplete or the institution is secretly in Regime 2 (has no true $S_0$).

\subsection{Cognitive and Neural Systems}

\textbf{Prediction 2.} A cognitive system that oscillates between clarity and confusion shows Regime 1 signature: periods of focus (near $S_0$) interrupted by necessary resets (sleep, attention shifts, off-loading to external memory). The reset cost (fatigue, cognitive cost of re-entry) should increase with the elapsed time since the last rest.

A cognitive system that cannot oscillate---that enters a state of permanent confusion, fragmentation, or dissociation---shows Regime 2 signature. Attempting to force it back to clarity is destabilizing (because there is no single $S_0$; there are multiple incompatible orderings, each partially correct). The system stabilizes in the fragmentation. Recovery requires resonance: finding a token (a sensation, a story, a practice) that connects all the incompatible frames.

\textbf{Falsification:} If a rested person (reset cost paid) does not show improved focus (lower $\xi$), then Regime 1 does not apply. If a fragmented person (Regime 2) recovers when forced to choose one ordering (suppressing all others), then the system never had Regime 2 incompleteness; it was always Regime 1 displaced, and suppression is the override that works for a while but accumulates return debt.

\subsection{Mathematical and Logical Systems}

\textbf{Prediction 3.} Formal systems (arithmetic, logic, computation, set theory) that show Regime 1 signature do not exist. No formal system oscillates toward and away from a ground state. All foundational systems show Regime 2 signature: they cannot decide their own fundamental questions, and the Wall is internal.

However, the theory predicts that resonance (Regime 3) is possible. The comparison isomorphisms of arithmetic geometry show this: two incompatible cohomological grammars sharing a token (the period integral, the motive). A proof of the Riemann Hypothesis via resonance would show this signature: the zeros of $\zeta(s)$ and the arithmetic primes sharing a frequency, both read by the same structure that neither grammar alone can construct.

\textbf{Falsification:} If a formal system is found that (1) oscillates between truth and falsehood, (2) has a ground state $S_0$ that is well-defined and internally recoverable, and (3) shows Regime 1 cost structure, then Strand A applies to formal systems. This would overturn the prediction that all fundamental systems live at the Wall. It would also be shocking.

\section{Conclusion}

The two strands of the corpus are one theory of holding versus letting go. Strand A describes the cost and frequency of return for systems that are displaced from a well-defined ground state. Strand B describes the mathematics of systems where the ground state is replaced by incompleteness, opposition, or a split aura---where the Wall is not a boundary but the system itself.

The three regimes are three distinct structures:
\begin{enumerate}
\item Regime 1: $S_0$ is well-defined, displacement is measurable, $C_{\text{return}}$ is finite. The system oscillates at $\tau^*$. Example: mycelia.
\item Regime 2: $S_0$ is missing or split, $C_{\text{return}}$ is undefined. The system lives in the Wall. Example: consciousness.
\item Regime 3: $S_0$ is replaced by a token, $C_{\text{return}}$ is local. The system returns to the frequency. Example: resonance.
\end{enumerate}

The Wall is not opposed to displacement. It is displacement at the singular limit, where the ground state becomes ill-defined and the return cost bifurcates. Othering is the structure of Regime 2 made explicit: the system that is defined by what it opposes, with no ground state to hold. Grammar (chimera, token, aura) is the mathematics of what must be held to keep the system coherent when the ground state is not a simple state but a balance, a tension, an interface.

The corpus is unified: one theory, two languages, three regimes, one question underneath.

How do you hold something that wants to break? And what breaks when you can no longer afford to hold?

\begin{thebibliography}{99}

\bibitem{maintenance} Rincón, D. (2024). Optimal Maintenance: The Calculus of Staying. Unpublished manuscript.

\bibitem{mycelia} Rincón, D. (2023). Mycelia: The Invisible Collective. Unpublished manuscript.

\bibitem{shinto} Rincón, D. (2023). Shinto as Maintenance: The Twenty-Year Rebuild. Unpublished manuscript.

\bibitem{sunrise} Rincón, D. (2023). Sunrise: The Daily Return. Unpublished manuscript.

\bibitem{consciousness} Rincón, D. (2024). Consciousness: The Land Space. Unpublished manuscript.

\bibitem{synthesis} Rincón, D. (2023). The Arithmetic Wall: A Synthesis. Unpublished manuscript.

\bibitem{five} Rincón, D. (2024). Five Walls, One Aura. Unpublished manuscript.

\bibitem{chimera} Rincón, D. (2024). Chimera: The Wall Made Living. Unpublished manuscript.

\bibitem{token} Rincón, D. (2024). The Token and the Blank. Unpublished manuscript.

\bibitem{resonance} Rincón, D. (2024). Resonance: When Two Grammars Read the Same Thing. Unpublished manuscript.

\bibitem{versus} Rincón, D. (2024). Versus as Attractor: Trump, Override, and the Politics Without a Ground State. Unpublished manuscript.

\bibitem{inhibition} Rincón, D. (2024). Inhibition: The Second Function of the Aura. Unpublished manuscript.

\end{thebibliography}

\end{document}
